\documentclass[conference]{IEEEtran}
\IEEEoverridecommandlockouts
\usepackage{cite}
\usepackage{amsmath,amssymb,amsfonts}
\usepackage{algorithmic}
\usepackage{graphicx}
\usepackage{textcomp}
\usepackage{xcolor}
\def\BibTeX{{\rm B\kern-.05em{\sc i\kern-.025em b}\kern-.08em
    T\kern-.1667em\lower.7ex\hbox{E}\kern-.125emX}}

\begin{document}

\title{LambdaScript : A Functional Programming Language}

\author{\IEEEauthorblockN{Alex Kozik}
\IEEEauthorblockA{\textit{Cornell University}\\
Ithaca, NY \\
ajk333@cornell.edu}
}

\maketitle

\begin{abstract}
We present LambdaScript, a statically-typed functional programming language inspired by Haskell and OCaml. LambdaScript features Hindley-Milner type inference, polymorphic algebraic data types, comprehensive pattern matching, and first-class functions with closures. The language supports advanced features including mutually recursive type definitions, exhaustiveness checking for pattern matches, list comprehensions, and operators as first-class values. This paper describes both the formal semantics and the complete implementation of LambdaScript, including the design of the type system, operational semantics, and the architecture of a full interpreter pipeline consisting of lexical analysis, parser combinators, constraint-based type inference, and environment-based evaluation. The implementation is validated through over 700 unit tests covering complex real-world algorithms including red-black trees, interpreters, and graph traversal.
\end{abstract}

\end{document}
